\chapter{Introduction}
\label{ch:intro}

\section{Context}

This report presents the final-year project in Computer Science Engineering at \ubi,
dedicated to the development of \rat: an integrated static and dynamic analysis platform
for detecting \emph{Remote Access Trojans} (RATs) in Windows executables (\texttt{.exe},
\texttt{.dll}) and, with reduced scope, C\# source files (\texttt{.cs}): on the web/API,
\texttt{.cs} files undergo string-based heuristics only (no PE imports or ILSpy); compilation
to \texttt{.exe} is available via the optional Tkinter GUI or manual build before upload.

RATs are among the most concerning classes of malware today: they allow an attacker to
remotely control a compromised system, harvest credentials, exfiltrate data, and execute
arbitrary commands. Their detection often requires combining \emph{offline} techniques
(static binary analysis) and \emph{online} techniques (behaviour observation in a controlled
environment), complemented by signature rules and scoring heuristics.

\rat\ was designed in an academic context as an \textbf{educational and triage} tool:
in addition to flagging risk indicators, it exposes decompiled code, structured reports, and
explanations of suspicious functions, enabling students and junior analysts to understand
\emph{why} a file is considered malicious. Complementarily, the web interface integrates a
\textbf{Google Gemini}-based assistant that allows users to ask, in natural language, about
the meaning of the pseudo-C excerpt visible in the results panel.

\section{Motivation}

The proliferation of malware with remote access capabilities demands accessible tools that
integrate multiple analysis layers into a single workflow. Commercial solutions such as
cloud sandboxes (VirusTotal, Any.Run, Hybrid Analysis) are powerful but often unsuitable for
teaching or local laboratory contexts, whether due to cost or sample confidentiality.

This project was motivated by the need to \textbf{integrate} static analysis, YARA rules,
decompilation, and sandbox execution in a single workflow spanning an API, a web interface, and
a desktop application. Equally important is \textbf{isolation}: potentially malicious samples
must run inside a Hyper-V virtual machine without Internet connectivity, with automatic restore
to a clean snapshot after each run. The platform must also favour \textbf{transparency}, producing
detailed reports---imports, C2 strings, evasion indicators, risk score, pseudo-C/IL---rather than
a binary verdict alone, with optional natural-language explanation of decompiled excerpts. Finally,
\textbf{academic reproducibility} requires open-source code, versioned documentation, automated
tests, and security guidance suitable for supervised laboratory deployment.

\section{Objectives}

The overarching goal of \rat\ is to deliver an integrated static and dynamic analysis platform
for RAT triage in an academic setting: the analyst uploads a suspect file, receives an
explainable risk assessment grounded in concrete indicators, and may optionally observe behaviour
in an isolated sandbox before deciding on further action.

To realise this goal, the project pursued several concrete aims. The system provides a Python
\textbf{backend} that receives samples, executes a complete static analysis pipeline, and
orchestrates dynamic analysis jobs in static, dynamic, or combined modes. A \textbf{web interface}
offers drag-and-drop upload, mode selection, progress streaming, report visualisation, decompiled
code panels, an optional Gemini assistant, and bilingual support (Portuguese/English). A
\textbf{WPF desktop application} serves as an alternative entry point with dependency bootstrap,
persistent language selection, and Hyper-V sandbox integration via Path~B. A \textbf{VM agent} and
PowerShell automation enable isolated sample execution and behavioural telemetry collection inside
the guest. An \textbf{explainable scoring model} (0--100) aggregates YARA matches, C2 indicators,
evasion techniques, packers, obfuscation, and persistence signals. Finally, the solution is
\textbf{validated} through automated tests, benign samples, and operational documentation for
local deployment.

\section{Document Organisation}

This report comprises six chapters, an appendix, and a bibliography. \textbf{Chapter~\ref{ch:estado-arte}}
(\emph{State of the Art}) surveys malware analysis techniques, commercial and open-source
alternatives, and positions \rat\ relative to existing sandboxes and detection approaches.
\textbf{Chapter~\ref{ch:tecnologias}} (\emph{Technologies and Tools Used}) describes the
multi-language stack---Python/FastAPI, React, .NET/WPF, YARA, Ghidra, ILSpy, Hyper-V, and
SQLite---together with security and isolation choices. \textbf{Chapter~\ref{ch:engenharia}}
(\emph{Software Engineering}) presents functional and non-functional requirements, UML diagrams,
the data model, trust zones, and the two sandbox orchestration paths. \textbf{Chapter~\ref{ch:implementacao}}
(\emph{Implementation}) details each software component, the user manual, experimental
evaluation, and the automated test suite. \textbf{Chapter~\ref{ch:conclusoes}} (\emph{Conclusions
and Future Work}) summarises achievements, acknowledges limitations, and outlines extensions such
as multi-VM scaling, Linux support, and MITRE ATT\&CK mapping. The \textbf{appendix} reproduces
representative source excerpts for the vm-agent and Hyper-V driver. The \textbf{bibliography}
lists the references cited throughout the document.
